%% kapa: exam=1 round=34 year=2023 period=1 form=A subject=real_estate subject_ko=부동산학원론 seq=18 number=98
%% source: Q-Net 감정평가사 기출문제 (cst00309) artlSeq=5211789 file=raw/1008/5211789/2023년도 제34회 감정평가사 1차시험 1교시.pdf
%% conversion: kordoc 4.24.51 -> tools/gen1.py (auto); review=proofread-sonnet; flags=none
\begin{question}{98}
부동산투자에 관한 설명으로 옳은 것을 모두 고른 것은?
\begin{framedtext}
ㄱ. 순현재가치(NPV)법이란 투자로부터 발생하는 현재와 미래의 모든 현금흐름을 적절한 할인율로 할인하여 현재가치로 환산하고 이를 통하여 투자의 사결정에 이용하는 기법이다. \par
ㄴ. 추계된 현금수지에 대한 위험을 평가하는 위험할증률의 추계는 투자기간의 결정 및 현금수지에 대한 예측 이전에 해야 한다. \par
ㄷ. 내부수익률(IRR)이란 투자로부터 발생하는 미래의 현금흐름의 순현재가치와 부동산가격을 1로 만드는 할인율을 말한다. \par
ㄹ. 수익성지수(PI)는 투자로 인해 발생하는 현금유입의 현재가치를 현금유출의 현재가치로 나눈 비율로서 1보다 크면 경제적 타당성이 있는 것으로 판단한다. \par
\end{framedtext}
\begin{choices}
  \item ㄱ, ㄹ
  \item ㄴ, ㄷ
  \item ㄱ, ㄴ, ㄷ
  \item ㄱ, ㄴ, ㄹ
  \item ㄱ, ㄴ, ㄷ, ㄹ
\end{choices}
\end{question}
\answer{1}
